\begin{resumo}
    A detecção precoce de lesões de pele é crucial para o tratamento de doenças graves, como o câncer de pele, mas o
    diagnóstico destas lesões requer um dermatologista. No Brasil, os \ac{acs} têm contato direto com a população,
    porém não possuem a qualificação necessária para realizar a triagem de lesões de pele. Diante disso, uma ferramenta
    que pudesse classificar essas lesões e fornecer pré-diagnósticos seria de grande auxílio. Neste contexto, os
    \acp{mllm} se mostram promissores, pois possuem a capacidade de interpretar imagens e gerar dados textuais.

    Este trabalho avaliou o uso do \ac{llama} 3.2 11B, para classificar lesões de pele e gerar laudos. Os experimentos
    foram divididos em duas etapas, com o \textit{fine-tuning} sendo realizado em conjuntos de dados diferentes em cada
    uma delas. Na primeira etapa, usando o conjunto de dados de dermatoscopia \ac{ham10000} apenas para classificação,
    o modelo alcançou uma acurácia de 86,2\% com \ac{lora}. Enquanto na segunda etapa, que visava a classificação e a
    geração de laudos com um conjunto de dados proveniente do \ac{sttsc}, o desempenho caiu significativamente, com o
    modelo atingindo uma acurácia de 45,2\% com \ac{qlora} e 44,5\% com \ac{lora}. O baixo desempenho dos modelos
    finais na classificação pode ser explicado por inconsistências no conjunto de dados.

    \textbf{Palavras-chave}: Lesões de pele. Câncer de pele. MLLM. LLaMA. PEFT. Fine tuning. QLoRA. LoRA.
\end{resumo}

\begin{abstract}
    \begin{otherlanguage*}{english}
        Early detection of skin lesions is crucial for treating serious diseases, such as skin cancer, but the
        diagnosis of these lesions requires a dermatologist. In Brazil, \acf{acs} have direct contact with the
        population, however, they do not have the necessary qualifications to perform skin lesion screenings.
        Therefore, a tool that could classify these lesions and provide pre-diagnoses would be of great help.
        In this context, \acfp{mllm} are promising, as they have the ability to interpret images and generate textual
        data.

        This work evaluated the use of the \acf{llama} 3.2 11B to classify skin lesions and generate reports. The
        experiments were divided into two stages, with fine-tuning performed on different datasets in each. In the
        first stage, using the \acf{ham10000} dermatoscopy dataset for classification only, the model achieved an
        accuracy of 86.2\% with \acf{lora}. In the second stage, which aimed at classification and report generation
        with a dataset from the \acf{sttsc}, the performance dropped significantly, with the model reaching an accuracy
        of 45.2\% with \ac{qlora} and 44.5\% with \ac{lora}. The low classification performance of the final models can
        be explained by inconsistencies in the dataset.

        \textbf{Keywords}: Skin lesions. Skin cancer. MLLM. LLaMA. PEFT. Fine-tuning. QLoRA. LoRA.
    \end{otherlanguage*}
\end{abstract}
